\section{Proofs of Some Lemmas in Preliminaries}\label{app: disequality reduction}
This appendix provides the deferred proofs of the preliminary lemmas.

\subsection{Proof of Some Lemmas in~\cref{subsec: tractable signatures}}\label{subapp: proof of subsec tractable signature}
This subappendix provides the deferred proof of the preliminary lemmas in~\cref{subsec: tractable signatures}.


\InvarianceUnderRealOrth*
\begin{proof}[Proof of~\cref{lem: invariance under real orthogonal transformation}]
    We prove the three statements separately.

    For the first statement, let $f\in \F$ has arity $n$. Since $O$ is real orthogonal matrix,
    we have 
    \begin{equation}
        \overline{Of}
        =\overline{O}^{\otimes n}\overline f 
        =O^{\otimes n}\overline f 
        =O\overline f.
    \end{equation}
    Therefore, if $\F$ is conjugate closed, then $\overline{f}\in \F$. Thus, $\overline{Of}=O\overline f\in O\F$. The converse follows by applying the same argument to the real orthgonal matrix $O^{-1}$. Hence, the first statement is hold.

    For the second statement, since holographic transformation preserves tensor decompositions into unary and binary signatures, let $f=f_1\otimes \cdots \otimes f_k$ where $f_i$ has arity at most two, then $Of=(Of_1)\otimes \cdots \otimes (Of_k)$.
    Applying the same argument to $O^{-1}$ gives
\[
    \mathcal F\subseteq\mathscr T
    \quad\Longleftrightarrow\quad
    O\mathcal F\subseteq\mathscr T.
\]
Now let
$
    \mathscr C\in\{\mathscr P,\mathscr A,\mathscr L\},
$
and suppose that $\mathcal F$ is $\mathscr C$-transformable, witnessed
by some $T\in\mathrm{GL}_2(\mathbb C)$. We claim that $TO^{-1}$
witnesses that $O\mathcal F$ is $\mathscr C$-transformable. First,
\[
    (TO^{-1})(O\mathcal F)
    =T\mathcal F
    \subseteq\mathscr C.
\]
Moreover, since $O$ is orthogonal, the binary equality signature is
invariant under $O$:
\[
    (=_2)O^{\otimes 2}=(=_2).
\]
It follows that
\[
\begin{aligned}
    (=_2)\bigl((TO^{-1})^{-1}\bigr)^{\otimes 2}
    &=(=_2)(OT^{-1})^{\otimes 2}\\
    &=\bigl((=_2)O^{\otimes 2}\bigr)
      (T^{-1})^{\otimes 2}\\
    &=(=_2)(T^{-1})^{\otimes 2}
      \in\mathscr C.
\end{aligned}
\]
Hence $O\mathcal F$ is $\mathscr C$-transformable. Applying the same
argument to $O^{-1}$ proves the converse. Therefore, all four
alternatives in condition \eqref{cond: tractable classes} are
invariant under $O$.

For the third statement, let $f$ be a nonzero signature. Suppose that $f=g\otimes h$, since the holographic transformation preserves tensor decompositions, we have $Of=Og\otimes Oh$. Thus, reducibility of $f$ implies reducibility of $Of$. Conversely, if $Of=g'\otimes h'$, then $f=(O^{-1}g')\otimes (O^{-1} h')$. Hence $f$ is reducible if and only if $Of$ is reducible, and $f$ is irreducible if and only if $Of$ is irreducible.
\end{proof}

\subsection{Proof of Some Lemmas in~\cref{subsec: hardness results}}
This subappendix provides the deferred proofs of the preliminary lemmas in~\cref{subsec: hardness results}
on tensor factorization and hardness from odd-arity signatures or
$\neq_4$, which are used in the orthogonality arguments and the
low-arity base cases of the main proof.

\CSPwithEQ*
\begin{proof}
    Let $\mathcal{G}=T(\F\cup \{=_2\})$.
    Suppose that $\#\operatorname{CSP}_2(T(\F\cup \{=_2\}))$ is not $\#\operatorname{P}$-hard, then by~\cite{real_holantc}, $\mathcal{G}\subseteq \mathscr{A},\mathscr{P},\mathscr{L}$ or $\mathcal{G}\subseteq \alpha\mathscr{A}$.

    We consider the first three cases. 
    
    If $\mathcal{G}\subseteq \mathscr{A}$, then $T\F\subseteq \mathscr{A}$ and $T(=_2)=TT^{\mathsf T}\in \mathscr{A}$. Let $E=TT^{\mathsf T}\in \mathscr{A}$. Then we can get $(=_2)T^{-1}=(T^{-1})^{\mathsf T}T^{-1}=(TT^{\mathsf T})^{-1}=E^{-1}$, thus $E$ is invertible. By~\cref{lm: affine singular value}, $E$'s two row has the same norm and either the two rows are proportional or orthogonal. The proportional property gives $\operatorname{rank}(E)=1$, which contradiction to the invertibility of $E$. Hence, $EE^\dagger=\rho I_2$ and $E^{-1}=\rho^{-1}E^\dagger$. Since $\mathcal{A}$ is closed under conjugation and transposition, $E^{-1}\in \mathscr{A}$. Therefore, $(=_2)T^{-1}\in \mathscr{A}$ and $\mathcal{F}\subseteq \mathscr{A}$-transformable. This contradicts the assumption that $\F$ does not satisfy condition~\eqref{cond: tractable classes}.

    If $\mathcal{G}\subseteq \mathscr{P}$, then $T\F\subseteq \mathscr{P}$ and $T(=_2)=TT^{\mathsf T}\in \mathscr{P}$. Let $E=TT^{\mathsf T}\in \mathscr{P}$. Then we can get $(=_2)T^{-1}=(T^{-1})^{\mathsf T}T^{-1}=(TT^{\mathsf T})^{-1}=E^{-1}$, thus $E$ is invertible. If the two variables of $E$ are independent, then $\operatorname{rank}(E)=1$ which contradicts to the invertibility of $E$. Thus, either $E$ is diagonal matrix or antidiagonal matrix, both give that $E^{-1}\in \mathscr{P}$. Therefore, $(=_2)T^{-1}\in \mathscr{P}$ and $\mathcal{F}\subseteq \mathscr{P}$-transformable. This contradicts the assumption that $\F$ does not satisfy condition~\eqref{cond: tractable classes}.

    If $\mathcal{G}\subseteq \mathscr{L}$, then $T\F\subseteq \mathscr{L}$ and $T(=_2)=TT^{\mathsf T}\in \mathscr{L}$. Let $E=TT^{\mathsf T}\in \mathscr{L}$. Then we can get $(=_2)T^{-1}=(T^{-1})^{\mathsf T}T^{-1}=(TT^{\mathsf T})^{-1}=E^{-1}$, thus $E$ is invertible. Let $\alpha=e^{\frac{\pi \mathfrak i}{4}}$ and $D:=\left( \begin{matrix}
        1 & 0\\
        0 & \alpha 
    \end{matrix}\right).$ The definition of $\mathscr{L}$ means that for any $(s_1,s_2)\in \mathscr{S}(E)$, we have $\left(\left( \begin{matrix}
        1 & 0\\
        0 & \alpha^{s_1}
    \end{matrix}\right)\otimes \left( \begin{matrix}
        1 & 0\\
        0 & \alpha^{s_2}
    \end{matrix}\right)\right)E\in \mathscr{A}$.

    If $00\in \mathscr{S}(E)$, then $E\in \mathscr{A}$. If $01\in \mathscr{S}(E)$, then $(I_2\otimes D)E\in \mathscr{A}$. If $10\in \mathscr{S}(E)$, then $(D\otimes I_2)E\in \mathscr{A}$. If $11\in \mathscr{S}(E)$, then $(D\otimes D)E\in \mathscr{A}$. This diagonal matrix does not change the support, thus the support of $E$ is also affine. If $|\mathscr{S}(E)|=1$, then contradicts to the invertibility of $E$. If $|\mathscr{S}(E)|=2$ and the nonzero position in the same row or column, then contradicts to invertibility. If $\mathscr{S}(E)=\{00,01,10,11\}$, then $E(x,y)=\lambda\mathfrak i^{Q(x,y)}$. And we have $E(x,y)/E(x',y')\in \{1,\mathfrak i,-1,\mathfrak{-i}\}$. In particular, $E(0,1)/E(0,0)=\mathfrak i^k$ for some $k$. Since $01\in \mathscr{S}(E)$, $(I_2\otimes D)E\in \mathscr{A}$ gives that $E'=\left( \begin{matrix}
        E(0,0) & \alpha E(0,1)\\
        E(1,0) & \alpha E(1,1)
    \end{matrix}\right)\in \mathscr{A}$. Thus, $\alpha E(0,1)/E(0,0)=\alpha \mathfrak i^k\notin \{1,\mathfrak i,-1,\mathfrak{-i}\}$. This contradicts to $E'\in \mathscr{A}$. If $\mathscr{S}(E)=\{01,10\}$, then we can suppose that $E=\left(\begin{matrix}
        0 & b\\
        b & 0
    \end{matrix}\right)$ since $E$ is symmetric. Since $E'=(I_2\otimes D)E=\left( \begin{matrix}
        0 & \alpha b\\
        b & 0
    \end{matrix}\right)\in \mathscr{A}$, but $\alpha b/b=\alpha\notin \{1,\mathfrak i,-1,\mathfrak{-i}\}$. This contradicts to $E'\in\mathscr{A}$.
    Lastly, $E=\left( \begin{matrix}
        a & 0\\
        0 & b
    \end{matrix}\right),$ and $b/a=\mathfrak i^k$. This gives $E=\lambda \operatorname{diag}(1,\mathfrak i^k)\in \mathscr{A}$. Also we can get $E^{-1}=\lambda^{-1}\operatorname{diag}(1,\mathfrak i^{-k})\in \mathscr{A}$.
    Finally, we can get that $\F\in \mathscr{L}$-transformable, which contradicts to $\F$ does not satisfies the condtion~\eqref{cond: tractable classes}.

    Next, we consider the case that $\mathcal{G}\subseteq \alpha \mathscr{A}$. Let $D:=\operatorname{diag}(1,\alpha)$, $E=TT^{\mathsf T}$ and $T'=D^{-1}T$. Since $T\F\subseteq \alpha \mathscr{A},T(=_2)\in \alpha \mathscr{A}$, we have $T'\F=D^{-1}T\F\in \mathscr{A}$ and $D^{-1}ED^{-1}\in \mathscr{A}$. Thus, after a holographic transformation $D^{-1}$, we get that $\#\operatorname{CSP}_2(D^{-1}T(\F\cup \{=_2\}))$ where $D^{-1}T(\F\cup \{=_2\})\in \mathscr{A}$, this case is equivalent to the first case $\mathcal{G}\in \mathscr{A}$.

    Eventually, $\#\operatorname{CSP}_2(T(\F\cup \{=_2\}))$ is $\#\operatorname{P}$-hard.
\end{proof}

\DisFourgivesHard*
\begin{proof}[Proof of~\cref{lm: disequality 4 gives dichotomy}]

    Let $g_4:=K^{\otimes 4}\neq_4$. We may discard zero signatures and handle nullary signatures as scalar factors. We divide the proof into three cases.

    \textbf{Case 1: $\F$ contains a nonzero signature of odd arity.}

    By the holographic transformation $K$, $\holant{\neq_2}{\neq_4,\hF}\equiv_T \holant{=_2}{g_4,\F}$, and by~\Cref{thm: holant odd} together with~\cite{GSS26}, $\holant{=_2}{g_4,\F}$ is either polynomial-time solvable or is $\#\operatorname{P}$-hard.

    \textbf{Case 2: every signature in $\hF$ is an $\eo$ signature.}
    Since $\neq_4$ is also an $\eo$ signature, then the problem $\holant{\neq_2}{\neq_4,\hF}$ is a $\ceo$ problem, namely $\holant{\neq_2}{\neq_4,\hF}\equiv_T \ceo(\neq_4,\hF)$. Then by~\Cref{thm: dichotomy for EO} together with~\cite{GSS26}, $\ceo(\neq_4,\hF)$ is either polynomial-time solvable or $\#\operatorname{P}$-hard.

    \textbf{Case 3: a non-$\eo$ signature occurs.}

    Let $\widehat{f}\in \hF$ have arity $2n$ and suppose that $\widehat{f}$ is not an $\eo$ signature. Then there exists $\alpha\in\mathrm{supp}(\widehat{f})$ such that $\mathrm{wt}(\alpha)=k\neq n$. We first show that, if $k<n$, we can realize a signature  $\widehat{g}$ of arity $r=2n-2k>0$ such that $\widehat{g}(\vec{0})\neq 0$, if $k>n$, we can realize a signature  $\widehat{g}$ of arity $r=2k-2n>0$ such that $\widehat{g}(\vec{1})\neq 0$. This follows from~\cite{Holant_Odd}.

    We treat the case $k<n$ firstly. Recall that we realized a signature $\widehat{g}$ of arity $r=2n-2k$ such that $\widehat{g}(\vec{0})\neq 0$. Let $a:=\widehat{g}(\vec{0})\neq 0$ and $b:=\widehat{g}(\vec{1})$. By the gadget construction in~\cite{cai-fu-shao-eo}, $\neq_{2r}$ can be realized on the RHS of $\holant{\neq_2}{\neq_4,\hF}$. We permute its variables so that $\mathrm{supp}(\neq_{2r})=\{0^r1^r,1^r0^r\}$. Then we connect the $r$ variables of $\widehat{g}$ with the first $r$ variables of $\neq_{2r}$ using $\neq_2$, leaving the last $r$ variables of $\neq_{2r}$ dangling. The resulting signature $\widehat{h}$ of arity $r$ is $\widehat{h}=[a,0,\dots,0,b]_r$. Since $\widehat{h}$ is realized on the RHS of $\holant{\neq_2}{\neq_4,\hF}$, we have $\holant{\neq_2}{\neq_4,\hF}\equiv_T \holant{\neq_2}{\neq_4,\widehat{h},\hF}$.

    We now distinguish whether $b$ is zero.
    
    \textbf{Subcase 3.1: $b=0$.}

    Then $\widehat{h}=[a,0,\dots,0]_r=a\Delta_0^{\otimes r}=(\sqrt[r]a \Delta_0)^{\otimes r}$. By~\Cref{lm: decomposition of tensor power}, we can get $\holant{\neq_2}{\neq_4,\hF}\equiv_T \holant{\neq_2}{\neq_4,\widehat{h},\hF}\equiv_T \holant{\neq_2}{\neq_4,\Delta_0,\hF}$. By a holographic transformation $K$, we can get $\holant{\neq_2}{\neq_4,\Delta_0,\hF}\equiv_T \holant{=_2}{g_4,u_0,\F}$ where $u_0=K\Delta_0=\frac{1}{\sqrt{2}}\begin{bmatrix}1\\ \mathfrak i\end{bmatrix}.$ Thus a nonzero odd arity signature is available, and by~\Cref{thm: holant odd} together with~\cite{GSS26}, $\holant{=_2}{g_4,u_0,\F}$ is either polynomial time solvable or $\#\operatorname{P}$-hard.

    \textbf{Subcase 3.2: $b\neq 0$.}

    Now $ab\neq 0$. We first suppose that $r\geq 4$. Let $q\in \mathbb{C}$ satisfying $q^{2r}=\frac{b}{a}$. Define 
    \[
        \widehat{Q}^{-1}
        =
        \begin{pmatrix}
        q^{-1}&0\\
        0&q
        \end{pmatrix},
        \qquad
        \widehat Q
        =
        \begin{pmatrix}
        q&0\\
        0&q^{-1}
        \end{pmatrix}.
    \]
    We have 
    \[
        (\widehat{Q}^{-1})^{\mathsf T}N_2\widehat Q^{-1}
        =
        \begin{pmatrix}
        q^{-1}&0\\
        0&q
        \end{pmatrix}
        \begin{pmatrix}
        0&1\\
        1&0
        \end{pmatrix}
        \begin{pmatrix}
        q^{-1}&0\\
        0&q
        \end{pmatrix}
        =N_2.
    \]
    By a holographic transformation $\widehat{Q}$,
    \[
    \begin{aligned}
        &\holant{\neq_2}{\neq_4,\widehat{h},\hF}\\
        &\equiv_T \holant{(\widehat{Q}^{-1})^{\mathsf T}N_2\widehat{Q}^{-1}}{\widehat{Q}\neq_4,\widehat{Q}\widehat{h},\widehat{Q}\hF}\\
        &\equiv_T \holant{\neq_2}{\widehat{Q}\neq_4,\widehat{Q}\widehat{h},\widehat{Q}\hF}
        \equiv_T \holant{\neq_2}{\neq_4,\widehat{Q}\widehat{h},\widehat{Q}\hF}.
    \end{aligned}
    \]
    Furthermore,
    \[
    \begin{aligned}
        \widehat{Q}\widehat{h}
        &=\widehat{Q}^{\otimes r}\widehat{h}
        =[aq^r,0,\dots,0,bq^{-r}]_r\\
        &=aq^r[1,0,\dots,0,1]_r=aq^r(=_r)
    \end{aligned}
    \]
    since $q^{2r}=b/a$. Thus $\holant{\neq_2}{\neq_4,\widehat{h},\hF}\equiv_T \holant{\neq_2}{\neq_4,=_r,\widehat{Q}\hF}$. By\cite{odd_holant_real}, we have
    \[
        \#\operatorname{CSP}_r(\neq_2,\neq_4,\widehat{Q}\hF)
        \equiv_T \holant{\neq_2}{\neq_4,=_r,\widehat{Q}\hF}.
    \]
    Then by~\cite{odd_holant_real}, $\#\operatorname{CSP}_r(\neq_2,\neq_4,\widehat{Q}\hF)$ is either polynomial-time solvable or $\#\operatorname{P}$-hard.

    It remains to consider $r=2$. In this case $\widehat{h}=[a,0,0,b]$. Recall that $\neq_8$ can be realized on the RHS of $\holant{\neq_2}{\neq_4,\widehat{h},\hF}$. Permuting the variables of $\neq_8$ so that $\mathrm{supp}(\neq_8)=\{0^41^4,1^40^4\}$.
    Connecting the first four variables of $\neq_8$ with two copies of $\widehat{h}$ using $\neq_2$, leaving the last four variables of $\neq_8$ dangling. The resulting quaternary signature is $\widehat{h}_4=[a^2,0,\dots,0,b^2]_4$. Since $a^2b^2\neq 0$, the preceding argument applies with $r=4$. This completes the case $k<n$.

    The case $k>n$ is similar.
\end{proof}


\ConjuDisFourHard*
\begin{proof}[Proof of~\cref{lem: In conjugate closed setting disequality 4 gives hardness}]
     By~\Cref{lm: disequality 4 gives dichotomy}, we just need to verify that $\holant{\neq_2}{\neq_4,\hF}$ will not fall into the polynomial-time solvable side of case 1, case 2, and case 3.

     \textbf{Case 1: $\F$ contains a nonzero signature of odd arity.}

     By~\Cref{lm: odd holant with conjugation}, this case is $\#\operatorname{P}$-hard.

     \textbf{Case 2: every signature in $\hF$ is an $\eo$ signature.}

     Suppose that $\hF\subseteq\POLUP$ and either $\hF\subseteq \eo^{\mathscr{A}}$ or $\hF\subseteq \eo^{\mathscr{P}}$. Since $\neq_4\in \POLUP\cap\eo^{\mathscr{A}}\cap\eo^{\mathscr{P}}$, then by~\Cref{lm: conjugation_polymorphism_affine support} we can get $\neq_4$ and every $\widehat{f}\in\hF$ has affine support. By~\cite{cai-fu-shao-eo}, these signatures are pairwise opposite. Combined with the definition of $\eo^{\mathscr{A}}$ and $\hF\subseteq \eo^{\mathscr{P}}$, we can get $\{\neq_4\}\cup \hF\subseteq \mathscr{A}$ or $\{\neq_4\}\cup\hF\subseteq\mathscr{P}$. Since $K^{\mathsf T}K=N_2\in\mathscr{A}\cap\mathscr{P}$, the transformation $K^{-1}$ witnesses $\mathscr{A}$- or $\mathscr{P}$-transformability of $\F$. This contradicts that $\F$ does not satisfy the condition \rm{(\ref{cond: tractable classes})}. The case $\hF\subseteq\POLDOWN$ and either $\hF\subseteq \eo^{\mathscr{A}}$ or $\hF\subseteq \eo^{\mathscr{P}}$ is symmetric. Thus this case is $\#\operatorname{P}$-hard.
     

     \textbf{Case 3: a non-$\eo$ signature occurs.}
     If exactly one of $a,b$ is nonzero in the preceding construction, we can get realize a nonzero unary signature, this case is $\#\operatorname{P}$-hard by~\Cref{lm: odd holant with conjugation}. Otherwise, by~\cite{odd_holant_real}, if $\F$ does not satisfy the condition \rm{(\ref{cond: tractable classes})}, $\#\operatorname{CSP}_r(\neq_2,\neq_4,\widehat{Q}\hF)$ is $\#\operatorname{P}$-hard for $r\geq 3$; when the original $r=2$, we first use the preceding quaternary construction and take $r=4$. Thus this case is also $\#\operatorname{P}$-hard.
\end{proof}











\subsection{Proof of Some Lemmas in~\cref{subsec: signatures and quantum states}}\label{subapp: proof of subsec signatures and quantum states}

\begin{proof}[Proof of~\cref{prop: invariant kth Orth}]
    For every subset $S\subseteq [n]$ of size $k$, if $f$ satisfies {\sc $k$th-Orth}, then there exists some $\mu\neq 0$ such that 
    \begin{equation}
        M_S(f)M_S(f)^\dagger=\mu I_{2^k}.
    \end{equation}

    We first prove invariance under conjugation. Since $M_S(\overline f)=\overline{M_S(f)}$, we have 
    \begin{equation}
        \begin{aligned}
            M_S(\overline f)M_S(\overline f)^\dagger
            &=\overline{M_S(f)}\,M_S(f)^{\mathsf T}\\
            &=\overline{M_S(f)M_S(f)^\dagger}\\
            &=\overline{\mu I_{2^k}}\\
            &=\mu I_{2^k}.
        \end{aligned}
    \end{equation}
   Therefore, $\overline{f}$ satisfies {\sc $k$th-Orth}. 

   For avoiding conflict between transpose $\mathsf{T}$ and holographic transformation $T$, we use $U$ instead of the holographic transformation $T$.
   We next prove invariance under uniatry holographic transformation $U\in \mathbf{U}_2(\mathbb{C})$. We know that $M_S(Uf)=U^{\otimes k}M_S(f)(U^{\mathsf{T}})^{\otimes(n-k)}$. Since $T$ is unitary, we have $U^{\mathsf{T}}\overline U=I_2$. Then we can get that 
   \begin{equation}
    \begin{aligned}
        M_S(Uf)M_S(Uf)^\dagger
        &= U^{\otimes k}M_S(f) (U^{\mathsf{T}})^{\otimes (n-k)}\overline{U^{\otimes (n-k)}}M_S(f)^\dagger (U^\dagger)^{\otimes k}\\
        &=U^{\otimes}M_S(f)M_S(f)^\dagger (U^\dagger)^{\otimes k}\\
        &=\mu U^{\otimes}I_{2^k} (U^\dagger)^{\otimes k}\\
        &=\mu I_{2^k}.
    \end{aligned}
   \end{equation}
   Thus, $Uf$ satisfies {\sc $k$th-Orth}.
\end{proof}
